\documentclass[10pt,letterpaper]{article}
\usepackage[utf8]{inputenc}
\usepackage[T1]{fontenc}
\usepackage[spanish]{babel}
\usepackage{amsmath, amssymb}
\usepackage[top=0.8cm, bottom=0.6cm, left=1.2cm, right=1.2cm, includeheadfoot]{geometry}
\usepackage{enumitem}
\usepackage{fancyhdr}
\usepackage{xcolor}
\usepackage{microtype}
\usepackage{array}
\usepackage{tabularx}
\usepackage{helvet}
\usepackage{tcolorbox}
\usepackage{graphicx}
\tcbuselibrary{skins,breakable}
\renewcommand{\familydefault}{\sfdefault}
\setlength{\headheight}{14pt}
\setlength{\parskip}{0pt}
\setlength{\parindent}{0pt}
\linespread{0.9}

\definecolor{applelightgray}{RGB}{180,180,180}

\pagestyle{fancy}
\fancyhf{}
\fancyhead[L]{\textbf{\sffamily Ejercicios de Pr\'actica \textendash\ Probabilidad y Estad\'istica}}
\fancyhead[R]{\small\sffamily Gu\'ia 5 \textendash\ Intervalos de Confianza}
\fancyfoot[C]{\small\sffamily\thepage}
\renewcommand{\headrulewidth}{0pt}
\fancyhead[C]{\textcolor{applelightgray}{\rule{\textwidth}{0.4pt}}}

\newcommand{\seccion}[1]{%
  \vspace{2pt}%
  \begin{tcolorbox}[enhanced,colback=white,colframe=black,arc=0pt,boxrule=0pt,
    leftrule=2pt,left=6pt,right=6pt,top=1pt,bottom=1pt]
    {\normalsize\bfseries\sffamily #1}
  \end{tcolorbox}%
  \vspace{1pt}%
}

\newcommand{\reglacaja}[1]{%
  \begin{tcolorbox}[enhanced,colback=white,colframe=black,arc=0pt,boxrule=0.5pt,
    left=6pt,right=6pt,top=2pt,bottom=2pt,breakable]
  \small\sffamily #1
  \end{tcolorbox}%
}

\begin{document}
\thispagestyle{empty}

\begin{center}
  {\fontsize{15}{17}\selectfont\bfseries\sffamily Gu\'ia 5 \textendash\ Intervalos de Confianza\par}
  \vspace{0.05cm}
  {\small\sffamily Estimaci\'on Puntual \textperiodcentered\ Margen de Error \textperiodcentered\ IC para Proporci\'on y Media\par}
  {\rule{3cm}{0.6pt}}
\end{center}

\vspace{0.15cm}

%% ===== PROBLEMA 1 =====
\seccion{Problema 1 \textendash\ IC para una Proporci\'on}

\noindent\sffamily\small Un centro de salud comunitario desea evaluar la proporci\'on de pacientes que consideran adecuada la atenci\'on m\'edica recibida. Se seleccion\'o una muestra aleatoria de $500$ pacientes; $415$ manifestaron estar conformes con la atenci\'on. Analice estad\'isticamente la proporci\'on de pacientes conformes usando una confianza del $95\%$:

\begin{enumerate}[leftmargin=1.2cm, label=\textbf{\alph*)}, itemsep=2pt, topsep=2pt]
  \item Determine la estimaci\'on puntual $\hat{p}$.
  \item Calcule el error est\'andar y el margen de error.
  \item Construya el intervalo de confianza del $95\%$ y concluya en contexto.
\end{enumerate}

\reglacaja{%
\textbf{Recuerde:} $\hat{p}=\dfrac{x}{n}$;\quad
$SE=\sqrt{\dfrac{\hat{p}(1-\hat{p})}{n}}$;\quad
$IC = \hat{p} \pm z_{\alpha/2}\cdot SE$, con $z_{0{,}025}=1{,}96$ para el $95\%$.%
}

%% ===== PROBLEMA 2 =====
\seccion{Problema 2 \textendash\ IC para una Media}

\noindent\sffamily\small En un estudio sobre frecuencia card\'iaca en reposo, se midi\'o a una muestra aleatoria de $36$ pacientes adultos, obteniendo un promedio de $72$ latidos por minuto con desviaci\'on est\'andar muestral de $9$ latidos.

\begin{enumerate}[leftmargin=1.2cm, label=\textbf{\alph*)}, itemsep=2pt, topsep=2pt]
  \item Construya un intervalo de confianza del $95\%$ para la frecuencia card\'iaca media de la poblaci\'on.
  \item Interprete el intervalo en contexto.
  \item Si se quisiera un intervalo m\'as angosto manteniendo la confianza, ?`qu\'e se debe hacer con el tama\~no muestral?
\end{enumerate}

%% ===== PROBLEMA 3 =====
\seccion{Problema 3 \textendash\ IC para Proporci\'on con Confianza del 90\%}

\noindent\sffamily\small En un programa de adherencia a tratamiento farmacol\'ogico, de una muestra aleatoria de $300$ pacientes cr\'onicos, $180$ declararon tomar sus medicamentos seg\'un lo indicado.

\begin{enumerate}[leftmargin=1.2cm, label=\textbf{\alph*)}, itemsep=2pt, topsep=2pt]
  \item Construya un intervalo de confianza del $90\%$ para la proporci\'on de adherencia ($z_{0{,}05}=1{,}645$).
  \item Interprete el intervalo en contexto.
  \item ?`Puede afirmarse con esta confianza que la adherencia supera el $50\%$? Justifique con el intervalo.
\end{enumerate}

\newpage

%% ================= RESPUESTAS =================
\seccion{Respuestas}

\small\sffamily
\textbf{Problema 1.}
\textbf{a)} $\hat{p}=\frac{415}{500}=\mathbf{0{,}83}$.
\textbf{b)} $SE=\sqrt{\frac{(0{,}83)(0{,}17)}{500}}\approx0{,}0168$; $ME=1{,}96\cdot0{,}0168\approx\mathbf{0{,}0329}$.
\textbf{c)} $IC_{95\%}=(0{,}7971\,;\,0{,}8629)$: con un $95\%$ de confianza, la verdadera proporci\'on de pacientes conformes est\'a entre $79{,}71\%$ y $86{,}29\%$.

\medskip
\textbf{Problema 2.}
\textbf{a)} $SE=\frac{9}{\sqrt{36}}=1{,}5$; $ME=1{,}96\cdot1{,}5=2{,}94$; $IC_{95\%}=(72-2{,}94\,;\,72+2{,}94)=(\mathbf{69{,}06\,;\,74{,}94})$ latidos/min.
\textbf{b)} Con un $95\%$ de confianza, la frecuencia card\'iaca media poblacional en reposo est\'a entre $69{,}06$ y $74{,}94$ latidos por minuto.
\textbf{c)} Aumentar $n$: el margen de error decrece proporcionalmente a $1/\sqrt{n}$.

\medskip
\textbf{Problema 3.}
\textbf{a)} $\hat{p}=\frac{180}{300}=0{,}60$; $SE=\sqrt{\frac{(0{,}6)(0{,}4)}{300}}\approx0{,}0283$; $ME=1{,}645\cdot0{,}0283\approx0{,}0465$; $IC_{90\%}=(\mathbf{0{,}5535\,;\,0{,}6465})$.
\textbf{b)} Con un $90\%$ de confianza, la proporci\'on de adherencia est\'a entre $55{,}35\%$ y $64{,}65\%$.
\textbf{c)} S\'i: todo el intervalo est\'a por sobre $0{,}50$, por lo que con esta confianza la adherencia supera el $50\%$.

\end{document}
